%%%%%%%%%%%%%%%%%%%%%%%%%%%%%%%%%%%%%%%%%%%%%%%%%%%%%%%%%%%%%%%%%%%%%%
\section{Architecture des données}
%%%%%%%%%%%%%%%%%%%%%%%%%%%%%%%%%%%%%%%%%%%%%%%%%%%%%%%%%%%%%%%%%%%%%%

La qualité d'un corpus dépend autant de sa structure que de la qualité des
enregistrements. Pour ALORA, l'organisation doit distinguer les fichiers
bruts, les données validées, les fichiers transcrits et les annotations,
afin de conserver une traçabilité complète de chaque étape.

L'organisation proposée garde une séparation nette entre le stockage des
fichiers audio et les métadonnées associées. Cela permet de préserver les
enregistrements originaux tout en rendant le corpus exploitable par des
outils de traitement automatique de la parole.

\begin{verbatim}
alora-corpus/
+-- audio/
|   +-- raw/
|   +-- validated/
|   +-- segmented/
+-- metadata/
|   +-- speakers.csv
|   +-- recordings.csv
|   +-- prompts.csv
+-- transcriptions/
|   +-- manual/
|   +-- reviewed/
+-- annotations/
|   +-- time/
|   +-- quality/
|   +-- linguistic/
+-- exports/
|   +-- train/
|   +-- validation/
|   +-- test/
+-- documentation/
\end{verbatim}

Les métadonnées principales concernent les locuteurs et les
enregistrements : identifiant, région, langue, variante, durée,
fréquence d'échantillonnage, statut de validation, date et prompt
associé. Les annotations, quant à elles, permettront d'enrichir chaque
fichier avec des informations temporelles, de qualité et linguistiques.

Au final, le corpus sera exporté dans un format compatible avec les
outils d'apprentissage automatique, avec une séparation explicite entre
les ensembles d'entraînement, de validation et de test.
